\section{Research Concept / Methodology}
\begin{enumerate}
    \item \textbf{Research Focus}
          \begin{itemize}
              \item Image focus detection as a classification problem:
                    \begin{itemize}
                        \item 2-class: focused vs unfocused
                        \item 3-class: focused, edge, unfocused
                    \end{itemize}
              \item Compare traditional machine learning methods vs lightweight neural networks
              \item Focus on low-sample scenarios
          \end{itemize}

    \item \textbf{Data Processing \& Segmentation}
          \begin{itemize}
              \item Divide images into patches / regions for classification
              \item Preprocessing: standardization, grayscale conversion
              \item Data augmentation for neural networks: rotation, flipping, noise addition
          \end{itemize}

    \item \textbf{Traditional Machine Learning Methods}
          \begin{itemize}
              \item \textbf{Algorithms:}
                    \begin{itemize}
                        \item Support Vector Machine (SVM)\cite{murphy2022pml}
                        \item Decision Tree\cite{murphy2022pml}
                        \item Random Forest\cite{murphy2022pml}
                        \item Optional: Clustering (for feature exploration / unsupervised analysis)
                    \end{itemize}
              \item \textbf{Features:}
                    \begin{itemize}
                        \item Edge strength, gradients, Laplacian
                        \item Texture descriptors (GLCM, LBP\cite{yi2016lbp})
                        \item Frequency-domain features (FFT, DCT)
                    \end{itemize}
              \item Advantages: fast training, interpretable
              \item Limitations: dependent on handcrafted features, limited robustness
          \end{itemize}

    \item \textbf{Neural Network Methods}
          \begin{itemize}
              \item \textbf{Algorithms:}
                    \begin{itemize}
                        \item Lightweight Convolutional Neural Network (CNN)\cite{zeng2018localmetric}
                        \item Optional: Transformer (attention mechanism for edge/blur focus)
                    \end{itemize}
              \item Characteristics:
                    \begin{itemize}
                        \item Automatic feature extraction
                        \item Potentially higher robustness
                        \item Sensitive to small sample sizes, mitigated by lightweight architecture \&
                              augmentation
                    \end{itemize}
          \end{itemize}

    \item \textbf{Experimental Design \& Comparison}
          \begin{itemize}
              \item Train both traditional and neural network models on the same dataset.
              \item Evaluate using multiple metrics:
                    \begin{itemize}
                        \item Classification accuracy (2-class / 3-class)
                        \item Precision, Recall, and F1-score
                        \item Confusion matrix
                        \item Robustness across different image types and blur levels
                        \item Training efficiency and data requirements
                    \end{itemize}
              \item Variables to investigate:
                    \begin{itemize}
                        \item Sample size
                        \item Classification type (binary vs multi-class)
                        \item Model type (traditional vs neural network)
                    \end{itemize}
          \end{itemize}

    \item \textbf{Expected Outcomes}
          \begin{itemize}
              \item Identify conditions where traditional ML is sufficient
              \item Determine scenarios where lightweight neural networks provide clear advantages
              \item Provide practical guidance for method selection in focus detection tasks
          \end{itemize}

\end{enumerate}